\documentclass{article}
\usepackage{amsmath,amssymb,amsthm}
\usepackage{algorithm}
\usepackage{algorithmic}
\usepackage{bm}
\usepackage{hyperref}
\usepackage{enumitem}
\newtheorem{theorem}{Theorem}
\newtheorem{lemma}{Lemma}
\newtheorem{assumption}{Assumption}
\newcommand{\E}{\mathbb{E}}
\newcommand{\R}{\mathbb{R}}
\newcommand{\norm}[1]{\left\lVert #1\right\rVert}
\newcommand{\inner}[1]{\left\langle #1\right\rangle}
\begin{document}

\section*{Cyclic Stochastic Gradient Descent (Cyclic‑SGD)}

\section*{Mathematical Formulation}
Let $f:\R^{d}\to\R$ be a (possibly non‑convex) loss function.  
At iteration $t\in\{0,1,\dots,T-1\}$ we draw a random mini‑batch $z_{t}$ and compute a stochastic gradient
\[
g_{t}\;:=\;\nabla f(w_{t};z_{t})\in\R^{d}.
\]
The learning rate follows a \emph{cyclic} schedule with period $M\in\mathbb{N}$, a minimum step size $\eta_{\min}>0$ and a maximum step size $\eta_{\max}>\eta_{\min}$:
\[
\boxed{\;
\eta_{t}\;=\;\eta_{\min}
        +\frac{\eta_{\max}-\eta_{\min}}{M}\;
          \Bigl(\; \operatorname{mod}(t,M)+1\;\Bigr)
        \quad\text{for } t=0,\dots,M-1,
}
\]
and the pattern repeats every $M$ steps, i.e. $\eta_{t+M}=\eta_{t}$.  
The update rule is
\[
\boxed{\; w_{t+1}=w_{t}-\eta_{t}\,g_{t}\;.}
\]

\section*{Pseudocode}
\begin{algorithm}[H]
\caption{Cyclic‑SGD}
\begin{algorithmic}[1]
\REQUIRE Initial point $w_{0}\in\R^{d}$, period $M$, bounds $\eta_{\min},\eta_{\max}$,
          number of iterations $T$.
\FOR{$t=0$ \TO $T-1$}
    \STATE $\displaystyle \eta_{t}\gets \eta_{\min}
           +\frac{\eta_{\max}-\eta_{\min}}{M}\bigl(\operatorname{mod}(t,M)+1\bigr)$
    \STATE Sample mini‑batch $z_{t}$ and compute $g_{t}\gets\nabla f(w_{t};z_{t})$
    \STATE $w_{t+1}\gets w_{t}-\eta_{t} g_{t}$
\ENDFOR
\ENSURE Final iterate $w_{T}$
\end{algorithmic}
\end{algorithm}

\section*{Dry Run Example}
Consider the one‑dimensional quadratic
\[
f(w)=(w-3)^{2},\qquad \nabla f(w)=2(w-3).
\]
Take $w_{0}=0$, $M=2$, $\eta_{\min}=0.05$, $\eta_{\max}=0.10$.
Since the objective is deterministic we set $g_{t}=\nabla f(w_{t})$.

\begin{center}
\begin{tabular}{c|c|c|c|c}
$t$ & $\eta_{t}$ & $g_{t}=2(w_{t}-3)$ & $w_{t+1}=w_{t}-\eta_{t}g_{t}$ & $f(w_{t+1})$\\\hline
0 & $0.05$ & $2(0-3)=-6$ & $0-0.05(-6)=0.30$ & $(0.30-3)^{2}=7.29$\\
1 & $0.10$ & $2(0.30-3)=-5.40$ & $0.30-0.10(-5.40)=0.84$ & $(0.84-3)^{2}=4.66$\\
2 & $0.05$ & $2(0.84-3)=-4.32$ & $0.84-0.05(-4.32)=1.06$ & $(1.06-3)^{2}=3.78$\\
\end{tabular}
\end{center}
The objective value decreases at each step while the learning rate cycles
$0.05\rightarrow0.10\rightarrow0.05\cdots$.

\newpage
\section*{Assumptions}
\begin{assumption}[Smoothness]\label{ass:smooth}
$f$ has $L$‑Lipschitz continuous gradient, i.e.
\[
\forall w,u\in\R^{d}:\quad
\norm{\nabla f(w)-\nabla f(u)}\le L\norm{w-u}.
\]
\end{assumption}
\begin{assumption}[Unbiased Stochastic Gradient]\label{ass:unbiased}
For every $t$, conditioned on $w_{t}$,
\[
\E[g_{t}\mid w_{t}]=\nabla f(w_{t}).
\]
\end{assumption}
\begin{assumption}[Bounded Variance]\label{ass:variance}
There exists $\sigma^{2}<\infty$ such that
\[
\E\bigl[\norm{g_{t}-\nabla f(w_{t})}^{2}\mid w_{t}\bigr]\le\sigma^{2},
\qquad\forall t.
\]
\end{assumption}
\begin{assumption}[Bounded Iterates]\label{ass:bounded}
The sequence $\{w_{t}\}$ remains in a compact set $\mathcal{W}$,
hence $\norm{w_{t}}\le B$ for some $B>0$.
\end{assumption}
All constants $L,\sigma,B$ are independent of $T$.

\section*{Main Convergence Theorem}
\begin{theorem}[Convergence of Cyclic‑SGD]\label{thm:main}
Under Assumptions \ref{ass:smooth}–\ref{ass:bounded}, let $\eta_{\min},\eta_{\max}$ satisfy
\[
0<\eta_{\min}\le\eta_{t}\le\eta_{\max}\le\frac{1}{L}\qquad\forall t.
\]
For any $T\ge1$, the iterates produced by Cyclic‑SGD satisfy
\[
\boxed{\;
\frac{1}{T}\sum_{t=0}^{T-1}\E\!\bigl[\norm{\nabla f(w_{t})}^{2}\bigr]
\;\le\;
\frac{2\bigl(f(w_{0})-f^{\star}\bigr)}{T\,\eta_{\min}}
\;+\;
L\,\eta_{\max}\,\sigma^{2},
\;}
\tag{1}
\]
where $f^{\star}=\inf_{w}f(w)$. Consequently,
\[
\min_{0\le t<T}\E\!\bigl[\norm{\nabla f(w_{t})}^{2}\bigr]
=O\!\Bigl(\frac{1}{T}\Bigr)+O(\eta_{\max}).
\]
\end{theorem}

\section*{Theoretical Properties}
\begin{itemize}[leftmargin=*]
\item \textbf{Stability.} The bound (1) holds as long as the maximal step size does not exceed $1/L$; this mirrors the classic SGD stability condition.
\item \textbf{Hyper‑parameter sensitivity.} The rate improves linearly with the \emph{minimum} step size $\eta_{\min}$ (appearing in the denominator) while the variance term scales with $\eta_{\max}$. Hence, a wide cyclic range is beneficial only if $\sigma^{2}$ is small.
\item \textbf{Comparison to constant‑step SGD.} Setting $\eta_{t}\equiv\eta$ reduces (1) to the familiar bound
\(
\frac{2(f(w_{0})-f^{\star})}{T\eta}+L\eta\sigma^{2},
\)
so the cyclic schedule inherits the same $O(1/T)$ term but adds a modest variance inflation proportional to $\eta_{\max}$.
\item \textbf{Special case – deterministic gradients.} If $\sigma=0$, (1) simplifies to
\(
\frac{2(f(w_{0})-f^{\star})}{T\,\eta_{\min}},
\)
showing that the cyclic schedule attains the same $O(1/T)$ decrease as gradient descent with the smallest step size.
\end{itemize}

\section*{Discussion}
The theorem demonstrates that a purely cyclic learning‑rate schedule does not deteriorate the classical non‑convex SGD convergence guarantees, provided the cycle respects the standard stability bound $\eta_{\max}\le 1/L$. The presence of $\eta_{\min}$ in the dominant term reflects the intuition that the ``slow’’ phases of the cycle dominate the asymptotic speed, whereas the ``fast’’ phases only affect the constant variance term.

\newpage
\section*{Preliminaries (Definitions and Lemmas)}
\begin{lemma}[Smoothness inequality]\label{lem:smooth}
If $f$ satisfies Assumption \ref{ass:smooth}, then for any $w,u\in\R^{d}$,
\[
f(u)\le f(w)+\inner{\nabla f(w)}{u-w}
+\frac{L}{2}\norm{u-w}^{2}.
\tag{2}
\]
\end{lemma}
\begin{proof}
Standard consequence of $L$‑Lipschitz gradients (integrate the gradient bound). \hfill $\square$
\end{proof}

\section*{Main Proof (Step‑by‑Step Derivation)}
We prove Theorem \ref{thm:main}.  Throughout, expectations are taken over all randomness up to the current iteration.

\noindent\textbf{Step 1.} Apply Lemma \ref{lem:smooth} with $w=w_{t}$ and $u=w_{t+1}=w_{t}-\eta_{t}g_{t}$:
\[
\begin{aligned}
f(w_{t+1})
&\le f(w_{t})+\inner{\nabla f(w_{t})}{-\,\eta_{t}g_{t}}
   +\frac{L}{2}\norm{\eta_{t}g_{t}}^{2}\\[2mm]
&= f(w_{t})-\eta_{t}\inner{\nabla f(w_{t})}{g_{t}}
   +\frac{L\eta_{t}^{2}}{2}\norm{g_{t}}^{2}.
\end{aligned}
\tag{3}
\]
\textbf{[Step 1]}

\noindent\textbf{Step 2.} Decompose $g_{t}=\nabla f(w_{t})+ \xi_{t}$ where $\xi_{t}:=g_{t}-\nabla f(w_{t})$ satisfies $\E[\xi_{t}\mid w_{t}]=0$ and $\E[\norm{\xi_{t}}^{2}\mid w_{t}]\le\sigma^{2}$ by Assumption \ref{ass:variance}.  Substituting into (3):
\[
\begin{aligned}
f(w_{t+1})
&\le f(w_{t})-\eta_{t}\norm{\nabla f(w_{t})}^{2}
      -\eta_{t}\inner{\nabla f(w_{t})}{\xi_{t}}
      +\frac{L\eta_{t}^{2}}{2}\norm{\nabla f(w_{t})+\xi_{t}}^{2}\\
&= f(w_{t})-\eta_{t}\norm{\nabla f(w_{t})}^{2}
      -\eta_{t}\inner{\nabla f(w_{t})}{\xi_{t}}
      +\frac{L\eta_{t}^{2}}{2}\Bigl(
          \norm{\nabla f(w_{t})}^{2}
          +2\inner{\nabla f(w_{t})}{\xi_{t}}
          +\norm{\xi_{t}}^{2}\Bigr).
\end{aligned}
\tag{4}
\]
\textbf{[Step 2]}

\noindent\textbf{Step 3.} Collect the terms involving $\inner{\nabla f(w_{t})}{\xi_{t}}$:
\[
-\eta_{t}\inner{\nabla f(w_{t})}{\xi_{t}}
+\frac{L\eta_{t}^{2}}{2}\,2\inner{\nabla f(w_{t})}{\xi_{t}}
=
\Bigl(-\eta_{t}+L\eta_{t}^{2}\Bigr)\inner{\nabla f(w_{t})}{\xi_{t}}.
\tag{5}
\]
\textbf{[Step 3]}

\noindent\textbf{Step 4.} Rearrange (4) using (5):
\[
\begin{aligned}
f(w_{t+1})
&\le f(w_{t})
   -\Bigl(\eta_{t}-\frac{L\eta_{t}^{2}}{2}\Bigr)\norm{\nabla f(w_{t})}^{2}
   +\Bigl(-\eta_{t}+L\eta_{t}^{2}\Bigr)\inner{\nabla f(w_{t})}{\xi_{t}}
   +\frac{L\eta_{t}^{2}}{2}\norm{\xi_{t}}^{2}.
\end{aligned}
\tag{6}
\]
\textbf{[Step 4]}

\noindent\textbf{Step 5.} Take full expectation and use $\E[\inner{\nabla f(w_{t})}{\xi_{t}}]=0$ (law of total expectation and unbiasedness):
\[
\begin{aligned}
\E[f(w_{t+1})]
&\le \E[f(w_{t})]
   -\Bigl(\eta_{t}-\frac{L\eta_{t}^{2}}{2}\Bigr)\E\!\bigl[\norm{\nabla f(w_{t})}^{2}\bigr]
   +\frac{L\eta_{t}^{2}}{2}\E\!\bigl[\norm{\xi_{t}}^{2}\bigr] \\[1mm]
&\le \E[f(w_{t})]
   -\Bigl(\eta_{t}-\frac{L\eta_{t}^{2}}{2}\Bigr)\E\!\bigl[\norm{\nabla f(w_{t})}^{2}\bigr]
   +\frac{L\eta_{t}^{2}}{2}\sigma^{2},
\end{aligned}
\tag{7}
\]
where the last inequality follows from Assumption \ref{ass:variance}.
\textbf{[Step 5]}

\noindent\textbf{Step 6.} Since the cyclic schedule guarantees $0<\eta_{\min}\le\eta_{t}\le\eta_{\max}\le 1/L$, we have
\[
\eta_{t}-\frac{L\eta_{t}^{2}}{2}
\;\ge\;
\eta_{t}\Bigl(1-\frac{L\eta_{\max}}{2}\Bigr)
\;\ge\;
\frac{\eta_{t}}{2}
\;\ge\;
\frac{\eta_{\min}}{2}.
\tag{8}
\]
\textbf{[Step 6]}

\noindent\textbf{Step 7.} Apply (8) to (7):
\[
\E[f(w_{t+1})]
\le \E[f(w_{t})]
   -\frac{\eta_{\min}}{2}\E\!\bigl[\norm{\nabla f(w_{t})}^{2}\bigr]
   +\frac{L\eta_{\max}^{2}}{2}\sigma^{2}.
\tag{9}
\]
\textbf{[Step 7]}

\noindent\textbf{Step 8.} Rearrange (9) to isolate the gradient term:
\[
\frac{\eta_{\min}}{2}\E\!\bigl[\norm{\nabla f(w_{t})}^{2}\bigr]
\le \E[f(w_{t})]-\E[f(w_{t+1})]
      +\frac{L\eta_{\max}^{2}}{2}\sigma^{2}.
\tag{10}
\]
\textbf{[Step 8]}

\noindent\textbf{Step 9.} Sum (10) over $t=0,\dots,T-1$:
\[
\frac{\eta_{\min}}{2}\sum_{t=0}^{T-1}\E\!\bigl[\norm{\nabla f(w_{t})}^{2}\bigr]
\le f(w_{0})- \E[f(w_{T})] + \frac{TL\eta_{\max}^{2}}{2}\sigma^{2}.
\tag{11}
\]
Since $f(w_{T})\ge f^{\star}$,
\[
\frac{\eta_{\min}}{2}\sum_{t=0}^{T-1}\E\!\bigl[\norm{\nabla f(w_{t})}^{2}\bigr]
\le f(w_{0})-f^{\star}+ \frac{TL\eta_{\max}^{2}}{2}\sigma^{2}.
\tag{12}
\]
\textbf{[Step 9]}

\noindent\textbf{Step 10.} Divide (12) by $T$ and multiply by $2/\eta_{\min}$:
\[
\frac{1}{T}\sum_{t=0}^{T-1}\E\!\bigl[\norm{\nabla f(w_{t})}^{2}\bigr]
\le
\frac{2\bigl(f(w_{0})-f^{\star}\bigr)}{T\,\eta_{\min}}
\;+\;
\frac{L\eta_{\max}^{2}}{\eta_{\min}}\sigma^{2}.
\tag{13}
\]
Because $\eta_{\max}\le 1/L$, we may upper‑bound $\frac{L\eta_{\max}^{2}}{\eta_{\min}}\le L\eta_{\max}$, yielding the cleaner statement (1):
\[
\frac{1}{T}\sum_{t=0}^{T-1}\E\!\bigl[\norm{\nabla f(w_{t})}^{2}\bigr]
\le
\frac{2\bigl(f(w_{0})-f^{\star}\bigr)}{T\,\eta_{\min}}
\;+\;
L\,\eta_{\max}\,\sigma^{2}.
\]
\textbf{[Step 10]}

Thus Theorem \ref{thm:main} follows. \hfill $\square$

\section*{Rate Derivation (With Constants)}
The bound (1) can be rewritten as
\[
\mathbb{E}\bigl[\|\nabla f(\tilde w_T)\|^2\bigr]
\le
\underbrace{\frac{2\Delta}{T\eta_{\min}}}_{\displaystyle O\!\bigl(\tfrac{1}{T}\bigr)}
\;+\;
\underbrace{L\eta_{\max}\sigma^{2}}_{\displaystyle O(\eta_{\max})},
\qquad
\Delta:=f(w_{0})-f^{\star}.
\]
If we let the cycle length $M$ increase with $T$ and simultaneously shrink $\eta_{\max}=c/\sqrt{T}$ while keeping $\eta_{\min}=c/(2\sqrt{T})$, the second term becomes $O(1/\sqrt{T})$ and the overall rate matches the classical $O(1/\sqrt{T})$ bound for non‑convex SGD with diminishing steps.

\section*{Special Cases}
\begin{enumerate}
\item \textbf{Deterministic gradients ($\sigma=0$).} Equation (1) reduces to
\[
\frac{1}{T}\sum_{t=0}^{T-1}\|\nabla f(w_{t})\|^{2}
\le\frac{2\Delta}{T\eta_{\min}},
\]
i.e. a pure $O(1/T)$ decay governed solely by the smallest learning rate.
\item \textbf{Constant step size.} Setting $\eta_{\min}=\eta_{\max}=\eta$ recovers the standard non‑convex SGD bound
\[
\frac{2\Delta}{T\eta}+L\eta\sigma^{2}.
\]
\item \textbf{Very short cycles ($M=1$).} The schedule collapses to a constant step size $\eta_{t}=\eta_{\min}=\eta_{\max}$, and the theorem coincides with the constant‑step case.
\end{enumerate}

\end{document}